\section {The drying watermelon (Arithmetic)}

A watermelon weighs $10$ kg and is $99\%$ water.
It sits in the sun and some of the water evaporates, until it is only $98\%$ water.
How much does the watermelon weigh now?

\subsection{Solution}

It weighs $5$ kg, half of what it weighed before.

The watermelon is $99\%$ water, so the part that is not water is $1\%$ of $10$ kg, which is $0.1$ kg.
Only water evaporates, so the part that is not water still weighs $0.1$ kg, but now it is $2\%$ of the watermelon.
$2\%$ is $\frac{1}{50}$, so the watermelon weighs $50 \cdot 0.1 = 5$ kg.

The answer looks surprising because a change from $99\%$ to $98\%$ sounds small.
But the part that is not water went from $1\%$ to $2\%$ of the weight, which means that the weight was cut in half.
